\subsection{Compression identities on the minimal branch}

The main-text compression identities are developed in the dedicated section on compression identities and carrier exclusion. This appendix subsection collects additional arithmetic continuations and the compression chain summary.

\begin{remark}[Abelian package from occupancy and carrier norm]
On the canonical branch,
\[
\Oadm\gamma=48\cdot\frac56=40.
\]
Thus the characteristic coefficient $41$ is compressed to
\[
41=\Oadm\gamma+1=40+1,
\]
that is, fermionic occupancy trace plus one light Higgs doublet contribution.
Writing $b_Y:=5\bone/3$ for the usual hypercharge-normalized coefficient and using the block-index relation $k_B=\frac32 I_1(B)$ to define $I_1^{\mathrm{EW}}:=\frac{5}{24}\bone$ and $k_{\mathrm{EW}}:=\frac{5}{16}\bone$, one obtains the unified chain
\begin{equation}
\label{eq:41-chain}
48\,I_1^{\mathrm{EW}}
=32\,k_{\mathrm{EW}}
=10\,\bone
=\Oadm\gamma+1
=41.
\end{equation}
Equivalently,
\[
I_1^{\mathrm{EW}}=\frac{41}{48},
\qquad
k_{\mathrm{EW}}=\frac{41}{32},
\qquad
\bone=\frac{41}{10},
\qquad
b_Y=\frac{41}{6}.
\]
This is an algebraic signature rather than an accidental coincidence: the same integer $41$ appears simultaneously as electroweak block trace, E$_8$ exponent, UV hypercharge trace, and carrier occupancy times norm plus one.
\end{remark}

\begin{remark}[Rank lock and single-seed compression: main-text reference]
The rank-lock identity $2^{g-3}=g-1\Rightarrow g=5$ is proved as \Cref{prop:rank-lock} in the main text.
The single-seed compression and the Cabibbo inversion $\phiz=\frac12(1-\sqrt{1-4\lambda_C^2})$ with all derived inter-observable relations are given in \Cref{eq:cabibbo-inversion,eq:lambda-c-master}. The carrier exclusion theorems for wrong discrete worlds are in \Cref{thm:carrier-exclusion}.
\end{remark}

\subsection{Optional arithmetic continuations}

The next relations mix the main-text carrier with optional appendix-level E$_8$, inflation, or geometric notation. They are useful compression patterns, but their interpretation should remain conditional or conjectural.

\begin{remark}[60, 48, and \texorpdfstring{$5/4$}{5/4} on the optional ladder]
If the optional scale ladder is started at $D_1=60$, then
\[
\frac{D_1}{\Oadm}=\frac{60}{48}=\frac54=B\gamma.
\]
On the minimal branch one may also write the same number as
\[
\frac54=\frac{\gcar}{\gcar-1}\qquad (\gcar=5).
\]
This compresses the E$_8$ start index, the admissible chiral count, the mirror factor, and the two-defect coefficient to the same rational factor.
\end{remark}

\begin{remark}[Factorization of the two-defect coefficient]
From the main-text identity
\[
\delta_2=\frac54\,\deltatop^2
\qquad\text{and}\qquad
\deltatop=\Oadm\,\cthree^4
\]
one obtains
\[
\delta_2=2880\,\cthree^8.
\]
On the canonical branch this factor admits the exact arithmetic decompositions
\[
2880=\Dstart\,\Oadm=60\cdot48,
\qquad
2880=|R(E_8)|\,\dim\mathfrak g_{\SM}=240\cdot12.
\]
The arithmetic identity is exact; the interpretation of this coincidence is still appendix-level.
\end{remark}

\begin{conjecture}[E$_8$ root-count reading]
On the same optional branch,
\[
|R(E_8)|=248-8=240=5\cdot48=5\,\Oadm.
\]
This suggests that the E$_8$ root count may be read as a fivefold unfolding of the admissible chiral packet, although that interpretation is not used as a main-text theorem.
\end{conjecture}

\begin{remark}[Reactor seed of the optional scale ladder]
If one writes the optional ladder as
\[
\varphi_n=\phiz e^{-\gamma}\left(\frac{D_n}{D_1}\right)^{\kappa_{E8}},
\qquad
D_n=60-2n,
\]
then the start value is precisely the reactor-angle seed,
\[
\varphi_n=\sin^2\theta_{13}\left(\frac{D_n}{D_1}\right)^{\kappa_{E8}}.
\]
Thus the appendix-level ladder may be read as a discrete continuation of the same seed that controls $\theta_{13}$.
\end{remark}

\begin{remark}[Conditional inflation-side compression]
If one imports the v2.8 Starobinsky-side relations
\[
A_s=\frac{N^2}{24\pi^2}\left(\frac{M}{\Mpl}\right)^2,
\qquad
\frac{M}{\Mpl}=\sqrt{8\pi}\,\cthree^4,
\]
then on the minimal branch
\[
A_s=\frac{N^2}{3\pi}\cthree^8=\frac{N^2}{8640\pi}\,\delta_2,
\]
and also
\[
\frac{M}{\Mpl}=\frac{\sqrt{8\pi}}{\Oadm}\,\deltatop.
\]
This is recorded only as an appendix-level compression identity because inflation is outside the present main-text claim surface.
\end{remark}

\begin{remark}[Dyonic intercept and the $\alpha$ kernel]
From the dyonic appendix the local astrophysical $\beta$ function reads
\[
\beta_{\mathrm{BH}}(r)=\frac{Q_eQ_m}{256\pi^4 r^2}=\frac{\deltatop}{3}\frac{Q_eQ_m}{r^2}=16\cthree^4\frac{Q_eQ_m}{r^2}.
\]
This is a cross-link between the EHT intercept channel and the $\alpha$ kernel: the same topological coefficient $\deltatop=48\cthree^4$ that pushes $\alpha$ into the precision zone also controls the local dyonic $\beta$ function. The local and cosmological sectors are therefore not two loose motifs but the same coefficient in two projections.
\end{remark}

\begin{remark}[Conditional occupancy route to gravity]
If the relative Einstein term is of the form
\[
S_{\mathrm{EH}}^{\mathrm{rel}}
=
\frac{\chi^2}{192\pi^2}\tr_{\mathcal H_{\mathrm{ch}}}\mathbf 1
\int d^4x\sqrt{-g}\,R,
\]
and if the family bridge identifies
\[
\tr_{\mathcal H_{\mathrm{ch}}}\mathbf 1=\Oadm=48,
\]
then
\[
\Mpl^2=\frac{\chi_0^2}{2\pi^2}.
\]
This is the occupancy-to-gravity route alluded to in the main text.
\end{remark}

\begin{remark}[Compression chain]
The appendix-level compression pattern may therefore be summarized as
\[
Y \leftrightarrow \inv,
\qquad
\gamma=\frac56,
\qquad
B\gamma=\frac54,
\qquad
\frac{D_1}{\Oadm}=\frac54,
\qquad
|R(E_8)|=5\,\Oadm,
\]
\[
\bone=\frac{\Oadm\gamma+1}{10},
\qquad
\Omega_b+\beta_{\mathrm{rad}}=\phiz,
\qquad
A_s=\frac{N^2}{8640\pi}\,\delta_2
\quad
\text{(conditional v2.8 import)}.
\]
\end{remark}

\subsection{Conditional minimal-parameter picture}

The next package is stronger but also more conditional. It assumes that the carrier rank, family multiplicity, and seam normalization are all read as closed rather than only partially inherited.

\begin{principle}[Conditional minimal-parameter compression]
Assume the canonical branch
\[
\gcar:=\rank E=5,
\qquad
N_{\mathrm{fam}}=3,
\qquad
\cthree=\frac{1}{8\pi}.
\]
Then the numerical kernel may be rewritten almost entirely in terms of
\[
(\pi,\gcar,N_{\mathrm{fam}})=(\pi,5,3).
\]
Indeed, on this branch one has
\[
\gamma=\frac56=\frac{\gcar}{\gcar+1},
\qquad
\dim S^+=2^{\gcar-1},
\qquad
\Oadm=N_{\mathrm{fam}}2^{\gcar-1}=48,
\]
\[
\phibase=\frac{1-\gamma}{\pi}=\frac{1}{(\gcar+1)\pi},
\qquad
\deltatop=\Oadm\,\cthree^4,
\qquad
\delta_2=\Dstart\,\Oadm\,\cthree^8,
\]
with
\[
\Dstart=\gcar\,\dim\mathfrak g_{\SM}=5\cdot 12=60,
\qquad
\Oadm=(\gcar-1)\dim\mathfrak g_{\SM}=4\cdot 12=48.
\]
\end{principle}

\begin{remark}[Compressed dependency chain]
Under the same conditional closure, the appendix-level numerical chain reads
\[
(\widetilde M\to M,\Seam,\inv,\chi)
\Longrightarrow
E_3\oplus E_2
\Longrightarrow
Y \leftrightarrow \inv
\Longrightarrow
S(U(3)\times U(2))
\Longrightarrow
S^+,
\]
\[
\Longrightarrow
N_{\mathrm{fam}}=3
\Longrightarrow
\Oadm=48
\Longrightarrow
(\deltatop,\delta_2,\phiz)
\Longrightarrow
(\alpha,\lambda_C,\beta_{\mathrm{rad}},\Omega_b,\sin^2\theta_{13}),
\]
with the optional E$_8$, inflationary, and gravity-side sectors then read as downstream continuations of the same compressed core.
\end{remark}

\section{Relative APS and superconnection setup}
\label{app:aps}

The geometric completion uses APS-type relative data \cite{aps1975}. In practice this means:
\begin{itemize}[leftmargin=1.5em, topsep=3pt, itemsep=2pt]
    \item one works with a declared pair $(D,D_{\mathrm{ref}})$ rather than a single unanchored operator,
    \item boundary conditions on the seam $\Seam$ are part of the definition of the relative index,
    \item the $\eta$-term keeps track of the residual spectral asymmetry,
    \item the superconnection package is used in the same spirit as the spectral-action tradition \cite{chamseddineConnes1997}, but here always with an explicit reference subtraction.
\end{itemize}

The appendix is intentionally short. Its role is to make the common reference structure visible without pretending that the full heat-kernel computation has already been completed.

\section{Yukawa kernels and positivity lemmas}
\label{app:yukawa}

The transport program is based on the regularized kernel
\[
\mathcal Y_f^{(\varepsilon)}=(D_f^\dagger D_f+\varepsilon^2)^{-1}.
\]
Three remarks are important.
\begin{enumerate}[leftmargin=1.5em, topsep=3pt, itemsep=3pt]
    \item The regulator $\varepsilon$ is not optional: without it the kernel becomes ill-defined exactly where the admissibility question is most delicate.
    \item Positivity of the kernel is a transport-side selection principle, not yet a proved theorem of the hard kernel.
    \item Residual prefactors $\Lambda_{f,j}$ are bounded transport data that still require an operator-level closure. They must remain visible and must not be rhetorically hidden as if the entire mass map were already derived without remainder.
\end{enumerate}

\begin{remark}[Positive-kernel determinant route to strong CP]
If the physical mass map is written in the polar form
\[
M_f=U_{f,L}H_fU_{f,R}^\dagger,
\qquad
H_f>0,
\qquad
U_{f,L},U_{f,R}\in SU(3)_F,
\]
with $H_f$ induced by the positive transport kernel $\mathcal Y_f^{(\varepsilon)}$, then the induced determinant phase vanishes:
\[
\arg\det M_f=0.
\]
At the same time weak CP may survive through
\[
V_{\mathrm{CKM}}=U_{u,L}^\dagger U_{d,L}.
\]
Combined with the sheet-odd cancellation of $\theta_{\mathrm{sheet}}$ in the relative action, this is the algebraic core of the conditional strong-CP route summarized in the main text.
\end{remark}

\section{Optional record algebra, horizons, and prediction semantics}
\label{app:record}

This appendix keeps the record-level constructions available without letting them dominate the scope of explicit claims in the main paper.

\subsection{Minimal admissible vacuum}

At finite volume or with a declared cutoff, an admissible vacuum candidate is written as
\[
\rho_B
:=
\argmin_{\rho\in\Vadm}\tr(\rho\Hrel),
\]
provided the relative Hamiltonian is lower bounded on the admissible domain. In the absence of a full existence theorem this should be read as a minimizing-sequence statement, not as a solved cosmological theorem.

\subsection{Record growth principle}

The optional record algebra is
\[
\Arec(t):=\bigoplus_{\Delta_i>\Hhub(t)}\mathcal A_i,
\qquad
S_{\mathrm{rec}}(t)=\log\dim\Arec(t).
\]
As $\Hhub(t)$ decreases, more stable channels open. The intended arrow-of-time claim is therefore
\[
\frac{dS_{\mathrm{rec}}}{dt}>0,
\]
but only as a record-growth principle, not as a replacement for detailed nonequilibrium dynamics.

\subsection{Observer algebra and Born-type evaluation}

An observer is modeled as a stable pointer subalgebra
\[
\Aobs\subset\Arec.
\]
With orthogonal projectors $P_i\in\Aobs$ one writes the measurement channel
\[
\mathcal M_O(\rho)=\sum_i(P_i\rho P_i)\otimes |i\rangle\langle i|.
\]
In finite-dimensional sectors, positivity, additivity on orthogonal decompositions, and phase invariance lead to the standard Born-type assignment
\[
p_i=\tr(\rho P_i).
\]
The claim is therefore only a constrained evaluation principle on admissible sectors, not a new theorem about measurement in all settings.

\subsection{Horizons as conjectural inner seams}

If a horizon is modeled as an inner seam, the manuscript uses only an approximate factorization language,
\[
\mathcal H_{\mathrm{out}}\otimes\mathcal H_H\otimes\mathcal H_{\mathrm{in}},
\]
or, better, an algebraic split on an appropriate code subspace. This appendix intentionally avoids selling the factorization as exact fundamental structure.

\subsection{Prediction semantics}

The appendix-level prediction map is
\[
Q=(\rho_0,\mathcal O_Q,\Delta_Q,t,\mathcal C_Q,\mathcal R_Q),
\qquad
\Pred(Q)=\mathcal R_Q\!\left(\tr(\rho_tE_Q(\Delta)),\mathcal C_Q\right),
\]
with
\[
\rho_t=U(t)\rho_0U^\dagger(t).
\]
This is kept only as an optional formal extension.

\section{Extended benchmarks}
\label{app:benchmarks}

\subsection{Extended comparison rows and supplementary checks}

\begin{table}[H]
\centering
\scriptsize
\begin{tabularx}{\linewidth}{@{}>{\raggedright\arraybackslash}p{0.18\linewidth}>{\raggedleft\arraybackslash}p{0.17\linewidth}>{\raggedleft\arraybackslash}p{0.17\linewidth}>{\raggedright\arraybackslash}p{0.22\linewidth}>{\raggedright\arraybackslash}p{0.14\linewidth}@{}}
\toprule
Observable & TFPT value & Reference value & Role & Status \\
\midrule
$\delta_{\mathrm{CKM}}$ & \num{1.198} rad & \num{1.196} rad & unresolved transport comparison & appendix comparison row \\
$m_{\pi^0}$ & \num{134.979} MeV & \num{134.977} MeV & hadronic out-of-sample check & appendix comparison row \\
$m_{bb}$ & \num{0.0015} eV & $< \num{0.2}$ eV & neutrinoless double-beta interface & appendix interface \\
$\eta_B$ & \num{5.97e-10} & \num{6.13e-10} & leptogenesis-side interface & appendix interface \\
$m_a$ & \num{65.2} $\mu$eV & direct-search target & axion mass interface & appendix interface \\
$\nu_a$ & \num{15.76} GHz & direct-search band & haloscope target frequency & appendix interface \\
\bottomrule
\end{tabularx}
\caption{Extended comparison rows kept out of the main text.}
\label{tab:extended-benchmarks}
\end{table}

